\documentclass[11pt]{article}
\usepackage[letterpaper,margin=1in]{geometry}
\usepackage[T1]{fontenc}
\usepackage{lmodern,amsmath,amssymb,amsthm,microtype}
\usepackage{hyperref}
\hypersetup{colorlinks=true,linkcolor=blue,citecolor=blue,urlcolor=blue,
  pdftitle={Canonical distinguishing colorings of cyclic tournaments},pdfauthor={Alex Chan}}
\newtheorem{theorem}{Theorem}[section]
\newtheorem{lemma}[theorem]{Lemma}
\newtheorem{corollary}[theorem]{Corollary}
\newcommand{\Z}{\mathbb Z}
\newcommand{\F}{\mathbb F}
\newcommand{\Sym}{\operatorname{Sym}}
\newcommand{\Aut}{\operatorname{Aut}}
\newcommand{\AGL}{\operatorname{AGL}}
\title{Canonical distinguishing colorings of cyclic tournaments}
\author{Alex Chan\\
\small Harvard University and NBER\\
\small \texttt{achan@hbs.edu}}
\date{October 30, 2026}
\begin{document}
\maketitle
\begin{abstract}
Albertson and Collins conjectured that every cyclic tournament of order
$2q+1$ is distinguished by coloring $0,\ldots,q$ with one color and the
remaining vertices with a second color. We prove the stronger statement
that, for every odd-order permutation group on $\Z_n$ containing the
translation $x\mapsto x+1$, the set $\{0,\ldots,(n-1)/2\}$ has trivial
setwise stabilizer. The proof combines induction on block systems with
the Burnside--Schur theorem. In the primitive case, intersection sizes
of translates of the interval restrict an affine stabilizer to the
identity or a reflection. The group theorem and its tournament
corollary are also formally verified in Lean.
\end{abstract}
\noindent\textbf{Keywords.} Cyclic tournament; distinguishing coloring;
regular subset; permutation group.\\
\textbf{Mathematics Subject Classification (2020).} 05C20, 20B25.

\section{Introduction}
A coloring of a digraph is \emph{distinguishing} if the identity is its
only color-preserving automorphism. For an odd integer $n=2q+1\geq3$,
a \emph{cyclic tournament} $T(n;S)$ has vertex set $\Z_n$ and arc
$i\to j$ precisely when $j-i\in S$, where $S\subseteq\Z_n\setminus\{0\}$
contains exactly one member of each pair $\{d,-d\}$ of nonzero residues.
Its automorphism group contains the translation $\tau:x\mapsto x+1$.

Albertson and Collins~\cite{AC} proposed a prescribed distinguishing
two-coloring for every cyclic tournament. In the formulation of
Meslem and Sopena~\cite[Conjecture 2]{MS}, the first color class is
$I_n=\{0,\ldots,q\}$. The original formulation~\cite[Conjecture 2.4]{AC}
uses $\{0,\ldots,q-1\}$; the statements are equivalent because the
complement of this latter set is $I_n+q$, and translations are
automorphisms. Meslem and Sopena proved the prescribed-coloring
statement for several classes, including Paley tournaments, and whenever
one of the two canonically induced subtournaments is rigid.
As they explain, the existence of some distinguishing two-coloring
already follows from Gluck's regular-subset theorem~\cite{Gluck}.
The issue here is the specified interval, independent of $S$.

A subset is \emph{regular} for a permutation group when its setwise
stabilizer is trivial. We prove the following group-theoretic statement.

\begin{theorem}\label{thm:main}
Let $n$ be a positive odd integer, and let $G\leq\Sym(\Z_n)$ have odd
order and contain $\tau:x\mapsto x+1$. Then
$I_n=\{0,1,\ldots,(n-1)/2\}$ is regular for $G$.
\end{theorem}

\section{The affine interval lemma}
The following elementary observation treats all odd primes, including
characteristic three, uniformly.

\begin{lemma}\label{lem:affine}
Let $p=2q+1$ be an odd prime and let $I=\{0,\ldots,q\}\subseteq\F_p$.
The setwise stabilizer of $I$ in $\AGL_1(p)$ consists precisely of the
identity and the reflection $x\mapsto q-x$.
\end{lemma}
\begin{proof}
For $d\in\F_p$, put $C(d)=|I\cap(I+d)|$. Taking $d$ to be its
representative in $\{1,\ldots,2q\}$ gives
\[
 C(d)=q+1-\min\{d,p-d\}.
\]
Indeed, for $1\leq d\leq q$, the intersection is $\{d,\ldots,q\}$;
the remaining values follow from $C(-d)=C(d)$. In particular, among
nonzero $d$ the largest value is $q$, attained exactly at $d=1,-1$.

Suppose $f(x)=\alpha x+\beta$, with $\alpha\neq0$, preserves $I$.
Since $f(I+d)=I+\alpha d$, bijectivity gives $C(\alpha d)=C(d)$.
Taking $d=1$ forces $\alpha=1$ or $-1$. If $\alpha=1$ and
$\beta\neq0$, then $I+\beta=I$ would give $C(\beta)=|I|=q+1$,
contrary to the displayed formula. Thus $\beta=0$ in this case.
The map $r(x)=q-x$ preserves $I$. If $\alpha=-1$, then $r\circ f$
is a translation preserving $I$, hence is the identity by the preceding
case. Therefore $f=r$.
\end{proof}

\section{Proof of the group theorem}
\begin{proof}[Proof of Theorem~\ref{thm:main}]
We use strong induction on odd $n$. The case $n=1$ is immediate.
Two observations will be used throughout: a subset and its complement
have the same setwise stabilizer, and, since $G$ contains every
translation, regularity of $I_n$ implies regularity of each translate
$I_n+t$. For the latter assertion, conjugate a stabilizer by $\tau^t$.
The action of $G$ is transitive because it contains $\tau$.

Suppose first that $G$ is imprimitive. Choose a nontrivial invariant
block system and let $B$ be the block containing $0$. If $t\in B$,
then $t\in B\cap(B+t)$, so the block property gives $B+t=B$.
It follows that $B$ is an additive subgroup of $\Z_n$. All blocks
are translates of $B$. Hence for odd $a,b>1$ with $n=ab$ the blocks are
\[
 B_r=\{r+at:0\leq t<b\},\qquad 0\leq r<a.
\]
Write $a=2u+1$ and $b=2v+1$. Since $(n-1)/2=av+u$, we have
\begin{equation}\label{eq:blocks}
 I_n\cap B_r=
 \begin{cases}
 \{r+at:0\leq t\leq v\},&0\leq r\leq u,\\
 \{r+at:0\leq t<v\},&u<r<a.
 \end{cases}
\end{equation}

Let $g\in G$ preserve $I_n$. Its induced permutation on the blocks
preserves their intersection sizes with $I_n$. The blocks with the
larger size $v+1$ are indexed precisely by $I_a$. The induced group
on the blocks has odd order and contains $r\mapsto r+1$ on $\Z_a$.
The induction hypothesis therefore makes the induced permutation of
$g$ the identity. Thus $g$ fixes every $B_r$ setwise.

For each $r$, let $H_r$ be the permutation group induced on $B_r$ by
its setwise stabilizer in $G$. It has odd order and contains the
restriction of $\tau^a$. In the coordinate $r+at\leftrightarrow t$
on $\Z_b$, this restriction is $t\mapsto t+1$. If $r\leq u$, the
preserved subset in~\eqref{eq:blocks} is $I_b$. If $r>u$, its
complement is $\{v,\ldots,2v\}=I_b+v$. In either case, induction,
together with the complement and translation observations, makes the
restriction of $g$ to $B_r$ the identity. Hence $g$ is the identity.

It remains to treat the primitive case. The Burnside--Schur theorem
in its permutation-action form~\cite[opening theorem]{EP} states that a primitive
finite permutation group containing a regular cyclic subgroup is
either $2$-transitive or permutationally isomorphic to a subgroup of
$\AGL_1(p)$ in prime degree $p$. A $2$-transitive group of degree
$n\geq3$ has order divisible by $n(n-1)$, which is even. Consequently
$n=p$ is prime and the affine alternative holds.

We verify that the affine description applies in the original residue
coordinates. Let $f:\Z_p\to\F_p$ be a bijection conjugating $G$ into
$\AGL_1(p)$. The permutation $f\tau f^{-1}$ has order $p$. An affine
map with multiplier different from $1$ has a fixed point and is
conjugate to a multiplication, so its order divides $p-1$.
Thus $f\tau f^{-1}$ is translation by some $c\neq0$. It follows that
$f(j)=f(0)+jc$ for every $j\in\Z_p$. The affine renormalization
$z\mapsto(z-f(0))/c$ makes this bijection the original residue
identification. Every element of $G$ is therefore affine in those
coordinates.

Lemma~\ref{lem:affine} now says that an element of $G$ preserving
$I_p$ is either the identity or $x\mapsto q-x$. The latter is a
nonidentity involution and cannot belong to an odd-order group.
This completes the induction.
\end{proof}

\section{Cyclic tournaments}
\begin{corollary}\label{cor:AC}
For every cyclic tournament $T(2q+1;S)$ with $q\geq1$, the coloring
\[
 \lambda(i)=
 \begin{cases}1,&0\leq i\leq q,\\2,&q<i\leq2q\end{cases}
\]
is distinguishing.
\end{corollary}
\begin{proof}
The group $\Aut(T)$ has odd order. Otherwise Cauchy's theorem would
give a nonidentity involution, which swaps two vertices and therefore
reverses the unique arc joining them, a contradiction. Moreover,
$\Aut(T)$ contains $\tau$. Theorem~\ref{thm:main} implies that the
only automorphism preserving the first color class $I_{2q+1}$ is the
identity, as required.
\end{proof}

\section{Formal verification and reproducibility}
Theorem~\ref{thm:main} and Corollary~\ref{cor:AC} have been formally
verified in Lean. The proof is available in the
\href{https://github.com/alexgocardinal/albertson-collins/tree/2802921/lean}{project repository}
at commit \texttt{2802921}. The declarations
\texttt{canonicalHalfGroup\_proved} and
\texttt{tournament\_half\_regular}, in the namespace
\texttt{AlbertsonCollins}, express the group theorem and the tournament
conclusion, respectively. The latter assumes that rotation of the
specified cyclic labels preserves all arcs; it concludes that the
canonical half has trivial setwise stabilizer in the automorphism group.

The formal proof includes the block actions, strong induction, and the
interval-intersection argument. The Burnside--Schur step uses proof
terms from AxiomMath's formalization~\cite{AxiomMath}, with the source
revision and Apache-2.0 license recorded in the repository. Thus the
classification step is proved rather than assumed as an additional
hypothesis. The complete project was built with Lean 4.35.0-rc3 and
a pinned Mathlib revision. The final theorem axiom reports contain
only \texttt{propext}, \texttt{Classical.choice}, and
\texttt{Quot.sound}. The finite computational checks are not premises
of either formal theorem.

\medskip
\noindent\textbf{Data availability.} The results are proved analytically;
no research datasets are used. The Lean proof, dependency pins,
computational scripts, and graph examples are available in the
\href{https://github.com/alexgocardinal/albertson-collins}{project repository}.

\medskip
\noindent\textbf{AI assistance.} ChatGPT/Codex (OpenAI) assisted with
checking the supplied proof, developing the interval-intersection
argument, implementing finite checks, developing and debugging the
Lean formalization, and drafting the exposition.

\begin{thebibliography}{9}
\bibitem{AC}
M.~O. Albertson and K.~L. Collins,
A note on breaking the symmetries of tournaments,
\emph{Congressus Numerantium} \textbf{136} (1999), 129--131.

\bibitem{AxiomMath}
AxiomMath,
\emph{Classification of Finite Simple Groups in Lean},
Burnside--Schur modules, source revision
\href{https://github.com/AxiomMath/CFSG/tree/69cb972f587a3ecc035e092253167ff2e303c962/CFSG/Permutation/BurnsideSchur}{\texttt{69cb972f587a}},
2026. Apache-2.0.

\bibitem{EP}
S. Evdokimov and I. Ponomarenko,
A new look at the Burnside--Schur theorem,
\emph{Bulletin of the London Mathematical Society}
\textbf{37} (2005), no.~4, 535--546.
\href{https://doi.org/10.1112/S0024609305004340}{doi:10.1112/S0024609305004340}.

\bibitem{Gluck}
D. Gluck,
Trivial set-stabilizers in finite permutation groups,
\emph{Canadian Journal of Mathematics} \textbf{35} (1983), no.~1, 59--67.
\href{https://doi.org/10.4153/CJM-1983-005-2}{doi:10.4153/CJM-1983-005-2}.

\bibitem{MS}
K. Meslem and E. Sopena,
On the distinguishing number of cyclic tournaments: \mbox{Towards} the
Albertson--Collins Conjecture,
\emph{Discrete Applied Mathematics} \textbf{266} (2019), 219--236.
\href{https://doi.org/10.1016/j.dam.2018.10.021}{doi:10.1016/j.dam.2018.10.021}.
\end{thebibliography}
\end{document}

